\section{Related work and novelty boundary}
\label{sec:related}

\subsection{Mobile GUI agents and benchmarks}

AppAgent uses multimodal interaction and experience learned through exploration or demonstrations \cite{zhang2023appagent}. Mobile-Agent emphasizes vision-centric perception without relying on app XML, while Mobile-Agent-v2 separates planning, decision, and reflection and Mobile-Agent-v3 broadens the framework across GUI environments \cite{wang2024mobileagent,wang2024mobileagentv2,ye2025mobileagentv3}. AutoDroid combines dynamic app analysis, UI representation, memory, and query optimization for Android task automation \cite{wen2023autodroid}. ScreenAgent and OSWorld expand the design and evaluation space for visual computer control \cite{niu2024screenagent,xie2024osworld}.

Android in the Wild provides large-scale demonstrations; AndroidEnv provides an interactive reinforcement-learning platform; AndroidWorld supplies reproducible state setup and programmatic task checks; and SPA-Bench integrates multiple smartphone agents and metrics \cite{rawles2023aitw,toyama2021androidenv,rawles2024androidworld,chen2024spabench}. These works establish observation and action interfaces that \SH\ must support. They do not claim an OS primitive for using an OTP or credential while withholding it from the agent. Moreover, AndroidWorld's standard applications avoid login-dependent tasks, which is why the proposed evaluation adds purpose-built, open authentication and checkout flows rather than claiming existing benchmark coverage.

\subsection{Placeholders and inference-time privacy}

The closest prior direction is privacy-preserving transformation of agent-visible content. Anonymization-Enhanced Privacy Protection for Mobile GUI Agents detects sensitive content, replaces it with type-preserving placeholders, transforms multiple UI representations, and resolves certain local interactions through a proxy \cite{zhao2026anonymization}. This demonstrates that semantic placeholders can preserve task utility. \SH\ changes the trust boundary in two ways: the handle is minted at a source that already owns or receives the value, so protection of integrated sources does not depend on finding the value after rendering; and use is authorized against an OS-observed destination rather than merely replacing a placeholder at a requested index.

Text and prompt privacy systems address a neighboring boundary. Casper detects and removes sensitive prompt content locally \cite{chong2024casper}; PAPILLON coordinates local and remote models for privacy-conscious delegation \cite{siyan2025papillon}; PrivacyRestore removes privacy spans and transmits a restoration representation in a client-server inference design \cite{zeng2025privacyrestore}; and resource-aware work estimates when sanitization will harm utility \cite{carpentier2024preempting}. These techniques can reduce cloud disclosure, but a mobile agent still needs a trusted actuation path when the hidden value must be entered into an application. We therefore treat sanitization as complementary and do not cite it as evidence for destination confinement.

\subsection{Prompt injection, agent isolation, and information flow}

Indirect prompt injection is an architectural risk when an agent interprets untrusted data as instructions \cite{greshake2023indirect}. InjecAgent and AgentDojo provide structured tests of attacks and defenses for tool-integrated agents \cite{zhan2024injecagent,debenedetti2024agentdojo}. Fine-print and environmental injection studies extend this concern to visual and GUI environments \cite{chen2025fineprint,liao2024eia,chen2025aeia}; mobile-agent experiments show that third-party channels can participate in realistic adversarial workflows \cite{du2025mobileattacks}.

Architectural defenses seek guarantees outside the language model. CaMeL extracts control and data flow from a trusted query and uses capabilities at tool calls \cite{debenedetti2025camel}. Fides assigns confidentiality and integrity labels in a formal planner model \cite{costa2025fides}. Progent enforces privilege policies for agents \cite{shi2025progent}. IsolateGPT separates apps and their interactions, while AirGapAgent limits task context to reduce disclosure under context hijacking \cite{wu2025isolategpt,bagdasarian2024airgap}. \SH\ shares the premise that a vulnerable model should not be the final policy authority. Its distinct object is a value that must cross from an OS source into a concrete UI field without crossing the planner's observation boundary.

\subsection{Capabilities, mobile information flow, and Android mediation}

Capabilities are unforgeable tokens of authority rather than descriptive names alone \cite{dennis1966semantics}. Least privilege, complete mediation, fail-safe defaults, and economy of mechanism remain useful design criteria \cite{saltzer1975protection}. Capsicum demonstrates a practical capability model for compartmentalizing commodity UNIX applications \cite{watson2010capsicum}. \SH\ is not a general-purpose object-capability OS; it uses a narrow capability-shaped record for one actuation right, including explicit attenuation by destination, task, time, and use count.

Android security research has shown both platform mediation and its limits. Android Permissions Demystified measured the mapping between API behavior and permissions \cite{felt2011permission}. Fine-grained mandatory access-control work and SEApp illustrate the value and engineering cost of adding policies below applications \cite{bugiel2013flexible,rossi2021seapp}. TaintDroid tracks sensitive flows dynamically from sources to sinks, while DroidScope reconstructs OS and runtime semantics for analysis \cite{enck2010taintdroid,yan2012droidscope}. \SH\ differs from taint tracking: it aims to prevent the raw value from entering an untrusted agent-visible representation, then authorizes one explicit sink. Tainting remains useful as an evaluation oracle for detecting accidental flows.

\subsection{Authentication and bound tokens}

HOTP and TOTP standardize common one-time-password constructions \cite{rfc4226,rfc6238}; NIST guidance describes authenticator and lifecycle considerations \cite{nist80063b}. These standards define or govern codes, not how an accessibility agent can enter a code without seeing it. Android's SMS Retriever can route a verification message associated with an app, and OTP Autofill can mediate user-approved filling \cite{androidSmsRetriever,androidAutofillServices}. Credential Manager and Autofill are primarily app- and user-driven platform interfaces \cite{androidCredentialManager,androidAutofill}. \SH\ reuses their source knowledge where available but adds an agent action that never returns the value to the agent.

Audience restriction in OAuth prevents a token intended for one protected resource from being replayed at another, and DPoP adds proof-of-possession sender constraint \cite{rfc8707,rfc9449,rfc9700}. WebAuthn scopes public-key credentials to relying parties \cite{webauthn3}. These systems motivate narrow binding and explicit validation. Their network principals and protocols differ from a package/origin/field target inside Android; consequently the paper claims a design analogy, not protocol equivalence.

\subsection{Novelty statement}

Within the verified literature above, we did not find a system that simultaneously (i) mints an agent-facing opaque handle at an Android authoritative source before screen exposure; (ii) binds resolution to the current task and OS-observed UI destination; (iii) enforces a type-specific lifecycle; and (iv) closes both accessibility and pixel observation after entry. This is a bounded literature statement, not a claim of exhaustive priority. The closest mobile-placeholder work is newer than the original concept and materially narrows the novelty boundary; any submission must re-run the search and compare implementations and threat models in detail.
